% ============================================================================
%  Chapter 7 --- Implementation
% ============================================================================
\chapter{Implementation}
\label{ch:implementation}

\section{Introduction}

This chapter details the technical implementation of the AI Operations Agent platform: the technology stack, service implementations, Docker composition, database schema deployment, and the end-to-end execution walkthrough. Every component described here is operational and containerised.

\section{Technology Stack}

\begin{table}[H]
\centering
\caption{Technology stack.}
\label{tab:tech-stack}
\begin{tabularx}{\textwidth}{@{}llX@{}}
\toprule
\textbf{Component} & \textbf{Technology} & \textbf{Version / Notes} \\
\midrule
Language          & Python 3                & All services written in Python \\
Web framework     & FastAPI                 & API Gateway + AI Service \\
ML library        & scikit-learn            & 1.5.2 (pinned in requirements.txt) \\
Data processing   & NumPy, pandas           & Feature engineering + data transforms \\
Serialisation     & joblib                  & Model persistence \\
Database          & PostgreSQL              & 16 (official Docker image) \\
Object storage    & MinIO                   & S3-compatible (official Docker image) \\
Containerisation  & Docker + Docker Compose & Multi-service orchestration \\
HTTP client       & httpx (async-capable)   & Pipeline worker $\to$ AI service calls \\
DB client         & psycopg2                & PostgreSQL driver \\
S3 client         & boto3                   & MinIO/OBS interaction \\
\bottomrule
\end{tabularx}
\end{table}

\section{Project Structure}

The repository follows a monorepo structure with each service in its own directory:

\begin{lstlisting}[language=bash, caption={Repository structure.}, label={lst:repo-structure}]
telecom-cloud-intelligence/
  docker-compose.yml
  services/
    api-gateway/
      Dockerfile
      main.py               # FastAPI (7 endpoints)
      requirements.txt
    ai-service/
      Dockerfile
      main.py               # ML models + 4 endpoints
      requirements.txt
    pipeline-worker/
      Dockerfile
      requirements.txt
      worker/
        __main__.py          # 5-phase pipeline
  docs/
    db/schema.sql            # PostgreSQL DDL
    architecture/            # Architecture docs
    data-model/              # Data lake + ER docs
  report/                    # This LaTeX report
  infra/                     # Infrastructure configs
\end{lstlisting}

\section{Service Implementation Details}

\subsection{API Gateway (\texttt{services/api-gateway/main.py})}

The API Gateway is a stateless, read-only service. It connects to PostgreSQL and serves the latest results for each analytical endpoint. Key implementation decisions:

\begin{itemize}
  \item \textbf{No cache layer:} Results are read directly from PostgreSQL. The data volume is small enough that query latency is negligible.
  \item \textbf{Pagination:} All list endpoints accept a \texttt{limit} parameter with sensible bounds (\texttt{ge=1, le=200/500}).
  \item \textbf{Error handling:} HTTP exceptions are propagated cleanly; internal errors return 500 with a message.
\end{itemize}

\begin{lstlisting}[language=Python, caption={API Gateway --- SLA risk endpoint (excerpt).}]
@app.get("/sla-risk")
def sla_risk_latest():
    with get_conn() as conn:
        with conn.cursor(cursor_factory=RealDictCursor) as cur:
            cur.execute("""
                SELECT id, run_id, region, window_start, window_end,
                       score, explanation, model_name, model_version,
                       created_at
                FROM sla_risk_scores
                ORDER BY created_at DESC LIMIT 1;
            """)
            row = cur.fetchone()
            if not row:
                raise HTTPException(404, "No SLA risk scores")
            return row
\end{lstlisting}

\subsection{AI Service (\texttt{services/ai-service/main.py})}

The AI Service hosts the three ML models and exposes inference endpoints. Key implementation decisions:

\begin{itemize}
  \item \textbf{Train-once, serve-many:} Models are trained on first startup (if no persisted artifacts exist) and reloaded on subsequent restarts from \texttt{/app/models/*.joblib}.
  \item \textbf{Docker volume (\texttt{aimodels}):} Model artifacts survive container restarts.
  \item \textbf{Deterministic fallback:} If the pipeline worker calls \texttt{/infer/sla-risk} without features, a deterministic hash-based score is returned (clearly marked as non-model-based).
  \item \textbf{Anomaly score normalisation:} Raw IsolationForest \texttt{decision\_function} scores are normalised to $[0, 1]$ for severity representation.
\end{itemize}

\begin{lstlisting}[language=Python, caption={AI Service --- IsolationForest inference (excerpt).}]
@app.post("/infer/anomaly")
def infer_anomaly(req: AnomalyRequest):
    X = np.array([
        [r.throughput_mbps, r.latency_ms, r.packet_loss_pct,
         float(r.active_users), r.signal_rsrp_dbm]
        for r in req.records
    ])
    labels = _anomaly_model.predict(X)  # +1 or -1
    raw_scores = _anomaly_model.decision_function(X)
    score_range = raw_scores.max() - raw_scores.min()
    norm_scores = 1.0 - (raw_scores - raw_scores.min()) \
                        / (score_range + 1e-9)
    # ... return anomaly results ...
\end{lstlisting}

\subsection{Pipeline Worker (\texttt{services/pipeline-worker/worker/\_\_main\_\_.py})}

The pipeline worker is a run-once container that executes the full 5-phase data pipeline. It orchestrates all other services:

\begin{enumerate}
  \item Creates MinIO buckets if they do not exist.
  \item Ingests data (real from \texttt{/data/tt-import/} or synthetic as fallback).
  \item Uploads raw data to MinIO \texttt{raw} bucket.
  \item Registers the pipeline run and datasets in PostgreSQL.
  \item Performs feature engineering (OSS: severity, QoS; BSS: ARPU category, intensity).
  \item Calls the AI Service for SLA risk inference.
  \item Calls the AI Service for OSS anomaly detection.
  \item Calls the AI Service for BSS revenue anomaly detection.
  \item Computes OSS--BSS statistical correlations.
  \item Builds the curated dataset and uploads to MinIO \texttt{curated} bucket.
  \item Persists all results to PostgreSQL.
  \item Registers model versions and marks the pipeline as succeeded.
\end{enumerate}

\section{Docker Compose Orchestration}

\begin{lstlisting}[language=yaml, caption={Docker Compose configuration (simplified).}]
services:
  postgres:
    image: postgres:16
    environment:
      POSTGRES_DB: telecom_intel
      POSTGRES_USER: telecom
      POSTGRES_PASSWORD: telecom_pw
    healthcheck:
      test: ["CMD-SHELL", "pg_isready -U telecom"]
    volumes:
      - pgdata:/var/lib/postgresql/data

  minio:
    image: minio/minio:latest
    command: server /data --console-address ":9001"
    healthcheck:
      test: ["CMD", "curl", "-f", 
             "http://localhost:9000/minio/health/live"]
    volumes:
      - miniodata:/data

  api-gateway:
    build: ./services/api-gateway
    ports: ["8000:8000"]
    depends_on:
      postgres: { condition: service_healthy }

  ai-service:
    build: ./services/ai-service
    ports: ["8001:8001"]
    volumes:
      - aimodels:/app/models
    depends_on:
      postgres: { condition: service_healthy }

  pipeline-worker:
    build: ./services/pipeline-worker
    depends_on:
      postgres: { condition: service_healthy }
      minio:    { condition: service_healthy }

volumes:
  pgdata:
  miniodata:
  aimodels:
\end{lstlisting}

\subsection{Health Checks and Dependencies}

Docker Compose health checks ensure services start in the correct order:
\begin{itemize}
  \item PostgreSQL must be healthy (\texttt{pg\_isready}) before API Gateway, AI Service, and Pipeline Worker start.
  \item MinIO must be healthy (HTTP liveness endpoint) before Pipeline Worker starts.
  \item The pipeline worker depends on both infrastructure services and the AI service (for inference calls).
\end{itemize}

\subsection{Named Volumes}

Three named volumes ensure data persistence across container restarts:
\begin{itemize}
  \item \texttt{pgdata} --- PostgreSQL data directory
  \item \texttt{miniodata} --- MinIO object storage
  \item \texttt{aimodels} --- Trained ML model artifacts (\texttt{.joblib} files)
\end{itemize}

\section{Database Schema Deployment}

The PostgreSQL schema is applied on first connection by the pipeline worker or via the \texttt{schema.sql} DDL file. Eight tables are created:

\begin{lstlisting}[language=SQL, caption={Key schema definitions (excerpt).}]
CREATE TABLE IF NOT EXISTS pipeline_runs (
  id         BIGSERIAL PRIMARY KEY,
  run_id     TEXT NOT NULL UNIQUE,
  status     TEXT NOT NULL DEFAULT 'started',
  started_at TIMESTAMPTZ NOT NULL DEFAULT now(),
  finished_at TIMESTAMPTZ,
  error_message TEXT
);

CREATE TABLE IF NOT EXISTS sla_risk_scores (
  id            BIGSERIAL PRIMARY KEY,
  run_id        TEXT REFERENCES pipeline_runs(run_id),
  region        TEXT,
  window_start  TIMESTAMPTZ,
  window_end    TIMESTAMPTZ,
  score         DOUBLE PRECISION CHECK (score >= 0 AND score <= 1),
  explanation   JSONB,
  model_name    TEXT DEFAULT 'sla-risk',
  model_version TEXT,
  created_at    TIMESTAMPTZ DEFAULT now()
);
\end{lstlisting}

\section{End-to-End Execution Walkthrough}

A complete platform execution follows these steps:

\begin{enumerate}
  \item \textbf{Start:} \texttt{docker compose up --build}
  \item PostgreSQL and MinIO start and become healthy.
  \item AI Service starts, loads/trains models, serves on port 8001.
  \item API Gateway starts, connects to PostgreSQL, serves on port 8000.
  \item Pipeline Worker starts:
  \begin{itemize}
    \item Checks for real TT data in \texttt{/data/tt-import/}
    \item If present: ingests, anonymises, maps to schema
    \item If absent: generates synthetic data (fallback)
    \item Executes 5-phase pipeline
  \end{itemize}
  \item \textbf{Query results:}
  \begin{itemize}
    \item \texttt{curl http://localhost:8000/sla-risk} --- latest SLA risk score
    \item \texttt{curl http://localhost:8000/anomalies} --- detected anomalies
    \item \texttt{curl http://localhost:8000/revenue-anomalies} --- BSS anomalies
    \item \texttt{curl http://localhost:8000/correlation} --- OSS--BSS correlations
  \end{itemize}
\end{enumerate}

\section{Chapter Summary}

This chapter detailed the technical implementation: Python/FastAPI services, scikit-learn models, Docker Compose orchestration, PostgreSQL schema, and the end-to-end execution flow. All services are fully containerised, use health checks for dependency management, and persist data across restarts through named volumes. The next chapter presents the evaluation results.
